% chktex-file 8   en-dash compounds are correct academic usage
% chktex-file 24  \label on own line is standard practice
% chktex-file 36  acronym parentheses (OUSD(A&S)) need no leading space
% chktex-file 38  American-style punctuation inside closing quotes is correct
\chapter{From Illumination to Assurance}
\label{chap:from_illumination_to_assurance}

This section translates the report's empirical findings into a policy agenda for reducing opacity in the Department of Defense semiconductor supply chain. Sections~\ref{chap:intro} through~\ref{chap:network_analysis} established the strategic stakes of semiconductor dependence, constructed a DoD-anchored multi-tier network from disclosed, predicted, and observed evidence, and identified high-consequence positions within that network using structural analysis and disruption testing. The task of Section~\ref{chap:from_illumination_to_assurance} is different. Rather than asking what the network looks like or which nodes matter most, this section asks what the Department of Defense and Congress should do with those findings.

The earlier sections show that it is possible to generate a decision-grade map of DoD-relevant semiconductor dependencies from unclassified and commercially available data. That map is auditable, reproducible, and useful for prioritization. It reveals hidden tiers, identifies chokepoints, and supports a structured workflow for deciding what to act on, what to validate, and what to monitor. At the same time, this report also demonstrates the limits of that achievement. The resulting network is not a complete census of the defense semiconductor supply chain. It is a defensible snapshot built from evidence streams that are themselves partial, uneven, and differently constrained. Residual opacity is not simply a flaw in the data. It is one of this report's core findings.

The policy problem that follows from this finding is clear. The Department of Defense cannot rely indefinitely on reverse-engineering its supply chain after the fact, whether through commercial illumination tools, episodic surveys, or one-off analytic efforts. Those approaches can improve visibility, but they do not alter the institutional conditions that generate opacity in the first place. Critical dependencies remain hidden because they are not systematically captured when suppliers are brought into a program, and they are not consistently updated as designs change, parts are substituted, ownership shifts, or production arrangements move across firms and borders. The issue, therefore, is not only how to see more. It is how to make the dependencies that matter more auditable across tiers and over time.

This section does not argue that the Department lacks regulations, priorities, or pilot authorities touching supplier assurance, program protection, lifecycle sustainment, counterfeit avoidance, or supply-chain monitoring. Existing policy already addresses mission-critical functions and trusted systems, lifecycle acquisition risk, diminishing manufacturing sources and material shortages (DMSMS), counterfeit electronic parts, and pilot supply-chain mapping efforts \parencite{ch5_dodi520044,ch5_dodi500085,ch5_tpp_guidebook_2022,ch5_dodi424515,ch5_dmsms_contract_language_guidebook,ch5_dfars2522467007,ch5_dfars2522467008,ch5_dfars2522047020,ch5_ndaa2024_sec856}. The argument is narrower and more consequential. The Department still lacks a coherent targeting and operating model for using those authorities to create decision-grade visibility for high-consequence semiconductor dependencies. The binding constraint is not whether the Department cares about supply chains. It is whether the Department can require the right evidence, from the right actors, at the right points in the lifecycle, and with the right validation priority \parencite{ch5_dod2022_supplychains,U.S.GovernmentAccountabilityOffice2025}.

The objective of this section is not perfect transparency. It is decision-grade visibility: a level of visibility sufficient to identify high-consequence dependencies, verify them when necessary, monitor meaningful change, and support timely intervention in the face of continuity and integrity risk. That standard is narrower than full supply-chain transparency but more demanding than retrospective illumination. It requires the Department to know enough, with sufficient confidence, to distinguish between what must be hardened immediately, what must be validated before action, and what can be monitored over time. The policy goal of this section is therefore fully continuous with the empirical goal of Section~\ref{chap:network_analysis}.

The core claim that follows from Sections~1--4 is therefore not that the Department needs another generic reporting rule. It is that the binding constraint is institutional evidence production. Illumination can prioritize where the Department should look, but assurance requires auditable objects: standardized, minimally disclosed evidence that can be required at specific acquisition touchpoints, refreshed when dependencies are rewired, and protected through controlled access. The recommendations in this section are therefore framed in terms of deliverables, triggers, responsibilities, and verification pathways rather than a mandated vendor or a single technical stack \parencite{ch5_nist800161r1upd1,ch5_dodi520044}.

To keep the recommendations realistic and coherent with the evidence presented thus far, this section focuses on acquisition, oversight, and traceability levers rather than prescribing a single technical architecture or arguing for universal disclosure of all supplier relationships. The emphasis is on policy requirements, governance principles, and institutional mechanisms that could reduce opacity while preserving legitimate protections for proprietary information and trade secrets. The focus remains defense-relevant semiconductor dependencies rather than the entire defense industrial base. The broader conclusion, including the project's overall limitations and future work, is reserved for Section~\ref{chap:conclusion}.

The remainder of Section~\ref{chap:from_illumination_to_assurance} proceeds in stages. Section~\ref{sec:ch5_decision_grade_illumination_limits} explains what Sections~1 through~4 accomplish and why a residual visibility ceiling remains. Section~\ref{sec:ch5_policy_objectives_design_principles} then formalizes the policy objective and design principles that should govern any credible effort to reduce supply-chain opacity under partial observability. Section~\ref{sec:ch5_layered_assurance_framework} translates those principles into a layered assurance framework that distinguishes traceability, provenance, resilience, and decision integration. The sections that follow operationalize this framework through acquisition and regulatory levers, data-governance choices, incentive and enforcement mechanisms, an implementation roadmap, and a set of specific recommendations for the Department of Defense and Congress. Throughout, this section uses the Section~\ref{chap:network_analysis} act-now, validate-then-act, and monitor workflow as the bridge between structural analysis and policy design.

%========================
% Section 5.1
%========================
\section{Decision-Grade Illumination and Its Limits}
\label{sec:ch5_decision_grade_illumination_limits}

This section is diagnostic. It summarizes what Sections~1--4 deliver for decision use and explains why a residual visibility ceiling remains even under a provenance-aware, decision-grade build. The section begins by identifying two core outputs that carry forward into policy design. The first is an auditable dependency map of DoD-relevant semiconductor relationships that records whether each tie is disclosed, predicted, or supported by observed shipping evidence. The second is a constrained decision workflow that translates structural rankings into an ``act now,'' ``validate then act,'' ``monitor,'' and ``deprioritize'' posture. This section also states what these outputs do not claim to provide. They are not a full census of the multi-tier supply chain; they do not trace specific components to specific defense platforms; and they do not directly test for integrity compromise. The visibility ceiling is then defined as the operational boundary produced by partial observability beyond Tier~1 and by limits in identity linkage at the demand anchor. The goal is to make the starting conditions for policy design explicit before Section~\ref{sec:ch5_policy_objectives_design_principles} formalizes the objectives and design principles that will govern credible recommendations. It also clarifies that the ceiling is not only a property of incomplete data. It is reinforced by a policy environment in which supplier visibility still sits across separate acquisition, sustainment, counterfeit, and program protection channels rather than a unified operating model \parencite{ch5_dod2022_supplychains,U.S.GovernmentAccountabilityOffice2025}.

Sections~\ref{chap:dod-supply-chain} and~\ref{chap:network_analysis} produce the main empirical output that this section builds on: a DoD-anchored map of who depends on whom in the semiconductor supply chain. The analysis starts from the firms the Department pays directly through prime contracts, then links those firms into the commercial identifier space to follow their upstream supplier relationships. The map is expanded using three kinds of evidence. Some links are disclosed in FactSet's Supply Chain Relationships data \parencite{factset_supply_2024}. Some links are added as high-confidence predictions from the Section~\ref{chap:link-prediction} models and are treated as hypotheses to be checked rather than facts. Some links are supported by observed shipping records, which provide transaction-style evidence that two firms exchanged goods at a point in time \parencite{factset_factset_2025}. For every link, the analysis retains a provenance tag that records which evidence stream supports it, so later analysis can distinguish what is directly evidenced from what is inferred. The result is auditable and reproducible, but it is still a defensible snapshot of DoD-relevant semiconductor dependencies rather than a complete census.

Section~\ref{chap:network_analysis} then uses this locked map to identify where continuity risk concentrates and what kinds of disruptions would matter most if they occurred. The analysis treats risk as a consequence problem rather than a forecasting problem. The analysis asks what happens to DoD-linked continuity when certain firms become unavailable or when paths through the network degrade, and compares targeted disruptions to a random-outage baseline. This stress-testing approach surfaces high-leverage positions such as chokepoints, bridges, and corridor intermediaries that sit on many dependency pathways. It also shows that some findings are robust across evidence views, while others shift when additional evidence layers are admitted. That stability check is part of what makes the output actionable because it separates results that can support immediate action from results that require follow-on validation work.

A ranking by itself is not yet a decision, so Section~\ref{chap:network_analysis} translates the stress-test results into a constrained workflow that matches how the Department actually operates. The workflow combines two pieces of information. One is impact, meaning how damaging disruption would be if a firm or pathway failed. The other is evidence stability, meaning whether a high-impact finding holds across alternative evidence-conditioned views of the same locked network. ``Act now'' covers findings that are high-impact and stable, so they justify immediate attention even under partial observability. ``Validate then act'' covers findings that appear high-impact but shift across views, so they warrant targeted supplier discovery and adjudication before committing resources. ``Monitor'' captures positions that are not immediate priorities but remain structurally consequential enough to watch for movement as the network changes. ``Deprioritize'' is an explicit constraint choice that prevents scarce analytic and engagement capacity from being spread too thin. This workflow is the section's bridge from structure to governance because it makes uncertainty operational, and it specifies what kind of follow-on work is required when visibility is incomplete.

This report's outputs are decision-useful in part because they draw hard boundaries around what the evidence can support. The map is not a complete census of the defense semiconductor supply chain, and some portions of the demand surface fall outside the mapped universe when prime recipients cannot be reconciled into commercial identifier space. Predicted links are not treated as ground truth, even when they are high-confidence, because their role is to focus validation effort rather than to certify dependency. The observed shipping layer strengthens confidence that two firms exchanged goods at a point in time, but it is still event-level evidence, and it does not trace specific components into specific defense platforms. For the same reason, Section~\ref{chap:network_analysis} is principally a continuity and structural severity section. It does not directly test tampering or compromise, and integrity enters here only as an implication of the continuity--integrity framework introduced in Section~\ref{chap:intro}. These boundaries prevent a category error in which illumination is mistaken for assurance.

The term ``visibility ceiling'' describes the operational boundary beyond which dependency evidence is not reliably produced, reconciled, and refreshed in a form that supports audit and decision use. One part of the ceiling is relational. Many upstream ties are never disclosed, are disclosed unevenly, or are disclosed without the standardization needed to integrate them cleanly across tiers. Another part of the ceiling concerns identity linkage. Even when a firm is known to DoD through contracting, it may not be linkable to the commercial identifier space that enables upstream mapping, creating structural blind spots at the demand anchor. The observed shipping layer helps by adding event evidence that two firms exchanged goods, but it does not cover the full universe of exchanges, and it is not designed to certify platform-specific provenance. These constraints are consistent with an institutional reality in which visibility is strongest near Tier~1 and degrades beyond it, and in which change over time can rewire dependencies faster than evidence is captured. Those institutional mechanisms are treated as explanations for why even a careful snapshot remains incomplete, not as outcomes directly measured in the data. That institutional reading is consistent with current assessments that DoD visibility efforts remain limited below the top tiers, unevenly coordinated, and not yet backed by a coherent set of tested reporting requirements across the enterprise \parencite{U.S.GovernmentAccountabilityOffice2025}.

The ceiling is most visible in practice through the structure of the Section~\ref{chap:network_analysis} triage outputs, especially the size and character of the ``validate then act'' queue. When a high-impact finding is stable across evidence views, the ceiling binds less, and the workflow can justify immediate action. When a high-impact finding shifts across views, the ceiling binds more and the correct response is targeted adjudication, not increased confidence. The validation queue is not a defect in the analysis. It is a diagnostic artifact of partial observability that identifies where the Department needs additional supplier discovery, documentation, or verification to move from plausible leverage to defensible action. The ``monitor'' category reflects the temporal side of the ceiling. Even when current impact is modest, limited change detection means the Department must watch for meaningful rewiring that can elevate risk. ``Deprioritize'' serves the same diagnostic purpose by making constraint explicit and by preventing scarce capacity from being consumed by low-yield leads.

Under the continuity lens introduced in Section~\ref{chap:intro} and operationalized through the disruption tests in Section~\ref{chap:network_analysis}, the policy relevance of the visibility ceiling is that it increases the probability of surprise propagation. When high-leverage dependencies sit beyond what can be evidenced, the Department is more likely to encounter deny and delay cascades only after they begin to disrupt schedules and readiness. This report shows that structural severity is concentrated and that conservative evidence conditioning can separate stable priorities from sensitive ones, but it also shows that uncertainty does not disappear. Existing policy already recognizes adjacent versions of this problem through program protection, counterfeit avoidance, DMSMS management, and pilot supply-chain mapping, but those authorities do not yet operate as a unified targeting model for semiconductor dependencies \parencite{ch5_dodi520044,ch5_dodi424515,ch5_dfars2522467007,ch5_dfars2522467008,ch5_ndaa2024_sec856}. The highest-value policy problem is therefore not how to produce a perfect map, but how to lower the ceiling where it matters most by making verification faster, cheaper, and more routine for high-consequence dependencies within an existing but still fragmented policy environment \parencite{ch5_dod2022_supplychains,U.S.GovernmentAccountabilityOffice2025}. The continuity--integrity framework from Section~\ref{chap:intro} implies a parallel constraint for integrity assurance. Event-level evidence and inferred links can accelerate who to investigate first, but they cannot certify item history, substitution accountability, or chain-of-custody into specific platforms. This diagnostic conclusion sets up Section~\ref{sec:ch5_policy_objectives_design_principles}, which formalizes objectives and design principles for credible interventions under partial observability.

%========================
% Section 5.2
%========================
\section{Policy Objectives and Design Principles}
\label{sec:ch5_policy_objectives_design_principles}

Section~\ref{sec:ch5_decision_grade_illumination_limits} makes the starting conditions for policy design explicit. This report produces an auditable, provenance-aware map and a constrained act now, validate then act, monitor, and deprioritize workflow, but it also shows that a residual visibility ceiling remains because upstream dependencies are only partially disclosed, only partially linkable, and only partially observable through event-level evidence. The policy task that follows is therefore not to promise perfect transparency or universal disclosure. It is to specify what kind of visibility is necessary for defensible action when continuity and integrity consequences are highest. This section formalizes that objective and then states the design principles that should govern any credible effort to reduce opacity in the DoD semiconductor supply chain under partial observability. The principles that follow do not ask the Department to care about an entirely new class of problem. They specify how existing authorities and existing priorities can be used in a more targeted, auditable, and lifecycle-persistent way for semiconductor-relevant dependencies \parencite{ch5_dodi520044,ch5_nist800161r1upd1,ch5_dod2022_supplychains}.

The objective is decision-grade visibility: enough timely, verifiable information about high-consequence dependencies to support defensible action without requiring a complete census of every firm, tier, and transaction. This standard is narrower than full transparency, but more demanding than retrospective illumination after disruption costs have already materialized \parencite{ch5_nist800161r1upd1}. In practice, it must enable the Department to identify which dependencies matter most, verify relationships when the evidence is thin or contested, detect meaningful change as designs and suppliers shift, and trace back quickly when disruptions, anomalies, or suspected compromise emerge. This objective is operationally defined by the workflow in Section~\ref{chap:network_analysis} because the purpose of visibility is not knowledge for its own sake, but the ability to act now when impact and confidence justify it, validate then act when leverage is plausible but stability is weak, monitor when positions are consequential but not urgent, and deprioritize when additional inquiry is unlikely to change a decision. To make this objective actionable, the policy response is organized around six design principles.

This visibility standard is also not a single requirement because the Department's exposure is not a single risk. As Section~\ref{chap:intro} argued, continuity risk concerns deny and delay: disruptions, bottlenecks, and outages that propagate into schedule failure, readiness degradation, or production shortfalls. Integrity risk concerns compromise: tampering, substitution, counterfeit insertion, or other forms of degradation that undermine trusted performance. These two risk modes create different evidence needs. Continuity decisions are often driven by dependency structure, substitutability, and change detection, while integrity decisions require auditable history, provenance, and substitution accountability. The design principles below are therefore intended to support both lenses, and Section~\ref{sec:ch5_layered_assurance_framework} makes the mapping explicit by distinguishing traceability, provenance, resilience, and decision integration as separate assurance functions \parencite{ch5_dodi520044,ch5_dod2022_supplychains}.

Principle~1 is risk-tiered scoping. Policy should require the deepest visibility and the strongest verification where continuity and integrity consequences are highest, and it should relax those demands as consequence falls. The visibility ceiling documented in Section~\ref{sec:ch5_decision_grade_illumination_limits} makes blanket disclosure and documentation requirements self-defeating because they either fail outright or generate uneven reporting that cannot support decision use. Operationally useful visibility therefore depends on selective depth that aligns evidentiary demands to the act now, validate then act, monitor, and deprioritize posture. It concentrates scarce analytic capacity, supplier engagement, and verification effort on the dependencies that plausibly drive disruption cascades. Risk-tiered scoping is therefore not a concession to opacity. It is a disciplined allocation of assurance effort \parencite{ch5_nist800161r1upd1}. It is a governance requirement for moving from one-off illumination toward durable assurance.

Principle~2 is verifiability over self-attestation. Section~\ref{sec:ch5_decision_grade_illumination_limits} showed that this report's map gains decision value by distinguishing among evidence types rather than treating all reported relationships as equivalent. Disclosed relationships are useful and often highly informative, especially for illumination, prioritization, and initial triage. But when a dependency carries serious continuity or integrity consequences, disclosure alone may not be enough to justify intervention, resolve ambiguity, or adjudicate meaningful change over time. That posture is consistent with Section~\ref{chap:link-prediction}, which treats predicted links as hypotheses to be checked rather than as facts to be assumed. The policy lesson is therefore not that reported relationships are untrustworthy per se. It is that high-consequence action should rest, wherever possible, on evidence that can be corroborated, reconciled, and audited when the stakes are highest \parencite{ch5_nist800161r1upd1,ch5_dodi520044}.

Principle~3 is minimum necessary disclosure. A credible visibility regime should require access to the information needed for decision use, but no more than that. The problem is not simply that too little information is publicly available. It is that the information required for continuity and integrity decisions is not reliably available to the right actors in a usable form. That does not require unrestricted transparency across the full multi-tier supply chain. Prime contractors may need visibility sufficient to manage and validate the dependencies that affect their own delivery obligations, while the Department may need a broader oversight view into high-consequence dependencies that cut across firms, programs, and tiers. Other commercial details may remain protected when they are not necessary for governance, verification, or intervention. Preserving trade secrets and legitimate proprietary interests is not in tension with this visibility objective \parencite{ch5_nist800161r1upd1,ch5_dodi520044}. It is part of the design condition for any regime that firms can participate in without treating disclosure itself as a source of strategic loss.

Principle~4 is interoperability and machine-readable standardization. Auditable supply-chain visibility cannot scale if each program, prime, or office records dependencies in its own format, with its own identifiers, and on its own update cycle. As Section~\ref{sec:ch5_decision_grade_illumination_limits} makes clear, visibility breaks down when dependency evidence cannot be reliably produced, reconciled, and refreshed for decision use, and Section~\ref{chap:dod-supply-chain} shows how much work is required even in a research setting to align disclosed, predicted, and observed evidence into a common provenance-aware snapshot. A policy regime that reproduces fragmentation at the point of reporting would simply recreate that same integration problem inside government. The design requirement is therefore not a single centralized database or a monolithic platform. It is a common, machine-readable structure for recording, exchanging, and refreshing dependency evidence so it can be reconciled across firms, programs, and tiers. At minimum, that requires a consistent identifier strategy, a small standardized set of fields sufficient for decision use, and a shared method for recording consequential change over time. Without that minimum standardization, the Department simply recreates inside government the same integration burden documented in Sections~\ref{chap:link-prediction} through~\ref{chap:network_analysis} \parencite{ch5_nist800161r1upd1}.

Principle~5 is lifecycle persistence and meaningful change detection. Supply-chain visibility cannot be treated as a one-time reporting exercise because redesign, obsolescence responses, supplier substitutions, ownership changes, geographic relocations, and ordinary supplier churn repeatedly rewire semiconductor supply chains. This point follows directly from the diagnostic conclusion of Section~\ref{sec:ch5_decision_grade_illumination_limits}, which shows that even a careful build remains a provenance-aware snapshot rather than a complete census, and from Section~\ref{chap:dod-supply-chain}, which treats the network as a locked snapshot so that Section~\ref{chap:network_analysis} can distinguish structural findings from ongoing change. A credible policy response must therefore ensure that visibility persists after initial disclosure and captures changes that materially alter continuity or integrity exposure \parencite{ch5_dod2022_supplychains,ch5_dodi424515}. The requirement is not constant surveillance of every transaction. It is a regime with update logic that makes consequential rewiring visible soon enough to support reclassification, revalidation, and intervention when the risk posture changes.

Principle~6 is burden-aware implementation under conditions of indirect reach. Visibility is strongest near Tier~1 and degrades beyond it, meaning the Department usually encounters sub-tier dependencies through prime contractors rather than direct contractual control. Any credible regime must therefore be designed for the structure through which information actually moves, not for an idealized system in which the government can interrogate every firm on demand. That constraint is especially important for smaller suppliers, which may face disproportionate reporting and verification burdens even when they occupy consequential positions in the network. A defensible visibility regime will not be durable if it depends on uniform demands that overwhelm supplier capacity or encourage defensive, low-quality reporting. A feasible regime must instead be phased, selective, and proportionate so participation remains administratively manageable, and the resulting information is more likely to be accurate, timely, and decision-useful \parencite{ch5_dod2022_supplychains,ch5_usc2223}.

Taken together, these principles define what a credible response to the visibility ceiling should seek to achieve and what constraints it must respect. The objective is decision-grade visibility, meaning enough timely and verifiable information about high-consequence dependencies to support defensible action under partial observability. The six principles specify how that objective should be pursued. Visibility should be risk-tiered, evidence should be verifiable, disclosure should be limited to what decision use requires, information should be structured for interoperability, visibility should persist through meaningful change, and implementation should remain feasible under conditions of indirect reach. These are not recommendations in themselves. They are design constraints for using existing authorities in a more coherent and targeted way across the rest of Section~\ref{chap:from_illumination_to_assurance}. The next section translates them into a layered assurance framework that distinguishes traceability, provenance, resilience, and decision integration so the Department can connect visibility requirements to practical governance choices without demanding perfect transparency.

%========================
% Section 5.3
%========================
\section{A Layered Assurance Framework}
\label{sec:ch5_layered_assurance_framework}

Section~\ref{sec:ch5_policy_objectives_design_principles} defined the objective of decision-grade visibility and set out the principles that should govern any credible effort to reduce opacity in the DoD semiconductor supply chain. Those principles establish what a workable policy response must achieve, but they do not yet specify how visibility should be organized in practice once the Department moves from diagnosis to institutional design. The next step is therefore to distinguish the different functions that an assurance regime must perform under partial observability. Some are about following dependencies and movements, some are about establishing evidentiary history, some are about understanding whether critical dependencies can withstand disruption, and some are about translating that information into defensible action. Those functions are organized as a layered assurance framework built around traceability, provenance, resilience, and decision integration. That distinction matters because current DoD practice already touches these evidentiary functions, but often through separate channels such as program protection, counterfeit avoidance, sustainment management, and emerging supply-chain mapping \parencite{ch5_dodi520044,ch5_dfars2522467007,ch5_dfars2522467008,ch5_dfars2522047020,ch5_dodi424515,ch5_ndaa2024_sec856}.

A layered framework is necessary because the visibility ceiling described in Section~\ref{sec:ch5_decision_grade_illumination_limits} is not a single informational deficit. The Department must be able to determine that a dependency exists, reconstruct where it originated and how it changed, assess whether disruption can be absorbed or recovered from, and decide what kind of intervention the evidence justifies. Those are related tasks, but they are not identical, and they do not require the same kind or degree of proof. A generic call for greater transparency therefore obscures the central policy problem: matching different decision questions to the forms of visibility and verification they actually require. The value of a layered framework is that it separates those functions without severing their connection, so the Department can ask not only how much it sees but also what kind of evidence it has and what that evidence is sufficient to support.

The framework is organized around four distinct but connected layers. Traceability concerns the structured visibility of dependencies, movements, custody changes, and other relevant supply-chain events across firms, tiers, and time. Provenance concerns the evidentiary history attached to those relationships and movements, including questions of origin, transformation, substitution, and documentation. Resilience concerns the evidence needed to judge whether a dependency can absorb disruption, adapt under stress, or recover after failure. Decision integration concerns how those three layers are translated into the Section~\ref{chap:network_analysis} act now, validate then act, monitor, and deprioritize workflow. The layers are therefore complementary rather than interchangeable, because each answers a different policy question and supports a different part of the Department's assurance problem.

Layer~1 is traceability. Traceability is the structured ability to follow relevant dependencies, movements, custody changes, substitutions, and other consequential supply chain events across firms, tiers, and time \parencite{ch5_dod2022_supplychains,ch5_nist800161r1upd1}. Sections~\ref{chap:dod-supply-chain} and~\ref{chap:network_analysis} show why this layer is foundational. Even under partial observability, the ability to locate where dependencies sit, how they connect, and how they may be rewired is what makes chokepoints, corridor intermediaries, and concentrated exposures visible for continuity analysis. In policy terms, traceability answers the first-order questions of who depends on whom, where a relevant relationship enters or exits the chain, and how a disruption might propagate across linked firms. What traceability does not do on its own is establish authenticity, resolve undocumented change, or determine whether what is being traced is sufficiently evidenced for adjudication.

Layer~2 is provenance. If traceability establishes that a dependency, component, or transaction moved through a chain, provenance establishes the evidentiary history attached to that movement. It asks where something originated, how it was transformed, who controlled it at consequential points, and whether undocumented substitution, alteration, or other material change occurred along the way. Provenance is therefore stronger than movement alone because it supports adjudication rather than simple observation. This layer matters most when the Department must assess authenticity, substitution, chain-of-custody, or other integrity-relevant questions that cannot be resolved by knowing only that a relationship exists or that an item changed hands. Provenance is therefore not a separate map of the supply chain. It is the evidentiary function that makes high-consequence dependencies more defensible to validate, contest, and act on when the stakes are highest \parencite{ch5_dodi520044,ch5_dfars2522467007,ch5_dfars2522467008}.

Layer~3 is resilience. In this framework, resilience refers to the evidence the Department needs to judge whether a high-consequence dependency can absorb disruption, adapt under stress, or recover after failure \parencite{ch5_dod2022_supplychains,ch5_dodi500085,ch5_dodi424515}. That includes questions such as whether there are qualified alternatives, whether substitution is feasible without unacceptable delay or degradation, whether surge or rerouting is possible, and whether the surrounding production and compliance conditions support continuity when a node is stressed. Resilience is therefore not a generic synonym for robustness. It is the layer that connects visibility to the practical question of what the Department can withstand, restore, or work around when a critical dependency is disrupted. If traceability shows where exposure sits and provenance shows what can be established about the dependency and its history, resilience shows whether that dependency is brittle, recoverable, or manageable under pressure. Resilience thus turns visibility from descriptive knowledge into evidence that can inform continuity planning and intervention.

Layer~4 is decision integration. Decision integration is the governance function that connects traceability, provenance, and resilience evidence to the workflow developed in Section~\ref{chap:network_analysis}. The network map identifies where the Department should look first by surfacing positions that are structurally consequential under partial observability. The layered assurance framework then helps determine what can be established about those positions, what remains uncertain, and what level of confidence is sufficient to support action. Decision integration does not replace structural analysis. It operationalizes it by linking evidentiary depth to the act now, validate then act, monitor, and deprioritize posture. Structural rankings prioritize attention, while assurance evidence helps adjudicate whether a dependency warrants immediate intervention, targeted validation, continued observation, or constraint-driven triage.

The four layers do not impose the same evidentiary burden on every dependency. Section~\ref{sec:ch5_policy_objectives_design_principles} argued that visibility must be risk-tiered, and that principle applies across the layered framework as well. Some dependencies may require only enough traceability to support monitoring and basic continuity awareness. Others may require stronger provenance to resolve ambiguity, and still others may require resilience evidence robust enough to judge whether substitution, rerouting, or intervention is actually feasible. The point is not to demand maximum depth everywhere. It is to match evidentiary depth to consequence so that the Department can escalate from visibility to verification and from verification to action where the stakes justify it. The layered framework is therefore scalable because it supports selective burden rather than uniform demand.

The layered framework is functional rather than technology-specific. Its purpose is to define what kinds of visibility and evidence the Department needs, not to prescribe a single platform, vendor, database, or reporting architecture. Section~\ref{sec:ch5_policy_objectives_design_principles} already established that credible policy design should privilege verifiability, interoperability, and burden-aware implementation, and those requirements can be met through more than one institutional or technical arrangement. What matters is that the evidence produced across the four layers can be reconciled, audited, refreshed, and used for decision-making when continuity and integrity consequences are serious. A technology-neutral stance is therefore a policy strength rather than a limitation because it keeps the framework focused on function, preserves flexibility across programs, and avoids mistaking one possible implementation path for the objective itself.

A layered assurance framework only matters if it persists through change. Section~\ref{sec:ch5_decision_grade_illumination_limits} showed that even a careful build remains a provenance-aware snapshot rather than a complete census, which means visibility degrades if redesign, obsolescence response, supplier substitution, ownership change, geographic movement, or other consequential rewiring is not captured over time. The value of the framework proposed here is therefore not only that it organizes what the Department should know, but also that it clarifies which changes must become visible early enough to trigger revalidation, renewed monitoring, or intervention. Traceability, provenance, resilience, and decision integration therefore provide a conceptual bridge from one-time illumination to durable assurance. The next section builds on that bridge by translating the layered framework into repeatable acquisition and oversight levers: specific touchpoints in the acquisition lifecycle where visibility must be created or refreshed, a small set of standardized evidence requirements that make visibility auditable rather than rhetorical, and a phased implementation logic tied to the Section~\ref{chap:network_analysis} act-now, validate-then-act, and monitor workflow through a more coherent use of existing authority channels.

%========================
% Section 5.4
%========================
\section{Acquisition and Regulatory Levers}
\label{sec:ch5_acquisition_regulatory_levers}

Section~\ref{sec:ch5_layered_assurance_framework} defined what kinds of evidence the Department needs to move from illumination to assurance. But a layered framework does not change practice unless it is embedded in the acquisition and oversight machinery that governs how programs are planned, awarded, executed, and sustained. The central claim of this section is therefore straightforward. The visibility ceiling persists not because the Department lacks concern, authority, or adjacent policy tools, but because critical dependency information is still not required, refreshed, and verified through a coherent operating model for high-consequence semiconductor dependencies \parencite{ch5_dod2022_supplychains,ch5_dodi500085,ch5_dodi520044,ch5_dodi424515,U.S.GovernmentAccountabilityOffice2025}. If structured semiconductor visibility is to become durable rather than episodic, it must enter the system as an institutional requirement rather than remain a retrospective analytic exercise.

Acquisition and oversight are the right policy channels because the Department does not control most semiconductor dependencies directly, especially below Tier~1, but it does shape contractor behavior through source selection, contractual deliverables, program protection requirements, reporting obligations, sustainment practices, quality controls, and audit or access rights \parencite{ch5_dodi500085,ch5_dfars2522467007}. The point is not to propose a wholly new bureaucracy before the Department has disciplined the tools it already possesses. It is to use existing acquisition machinery in a more targeted way for high-consequence semiconductor dependencies. The section therefore focuses on four linked design questions. It asks where visibility requirements should enter the lifecycle, what recurring artifacts must exist, how obligations should move beyond prime contractors, and how verification should be prioritized when analytic and oversight capacity is limited. Those questions convert the visibility objective from a general policy aspiration into an acquisition operating model.

Pre-award is where visibility must first become relevant for designated high-consequence acquisitions. The point at this stage is not to require universal multi-tier disclosure before contract award, which would be both unrealistic and poorly targeted. It is to make supplier visibility capacity, dependency identification processes, and assurance readiness part of how the Department evaluates risk in programs where semiconductor dependencies are likely to be mission-relevant, concentrated, or difficult to replace \parencite{ch5_dodi500085}. In practical terms, this means that for designated categories the Department should evaluate not only whether an offeror can deliver the end item, but also whether it has a credible process for identifying critical dependencies, producing and refreshing required visibility records, notifying the government when consequential changes occur, and supporting validation when asked. Visibility demanded only after award is often partial, expensive, and reactive. Pre-award is therefore the first point at which the Department can distinguish between contractors that merely promise performance and contractors that can support assurance.

The award stage is where expectation becomes obligation. If pre-award establishes that auditable supply-chain visibility matters for designated high-consequence acquisitions, award is the point at which the Department must convert that expectation into enforceable requirements that persist through performance \parencite{ch5_dodi500085,ch5_dfars2522467007,ch5_dfars2522467008}. For designated categories, contracts should not rely on broad statements that the contractor will manage supply-chain risk. They should require a small, standardized set of recurring visibility and assurance deliverables, define the conditions under which those records must be refreshed, and preserve the government's ability to review the evidence needed for validation. That is the point at which traceability, provenance, resilience, and decision use stop being policy aspirations and become part of program execution. Award is therefore the decisive institutional moment because it is where the Department can make visibility comparable across vendors, auditable across time, and actionable when high-consequence dependencies move from plausible concern to concrete risk.

A visibility regime only becomes real when it attaches to identifiable artifacts. For designated high-consequence dependencies, the Department should require a small set of recurring records rather than open-ended narrative reporting. At minimum, these should include a dependency disclosure package that identifies the critical supplier relationship or node at issue, a provenance evidence record that documents origin, custody, and relevant transformation or substitution history, a change notification that reports consequential rewiring such as supplier substitution, ownership change, geographic movement, or obsolescence workarounds, and a resilience support annex that states what qualified alternatives, fallback options, or recovery constraints bear on continuity planning \parencite{ch5_dodi424515,ch5_dfars2522467007,ch5_nist800161r1upd1}. These labels are less important than the underlying rule. Audit-ready visibility must attach to specific objects that can be furnished, reviewed, refreshed, and audited over time. Without that discipline, supplier visibility remains rhetorical because no one can tell what exactly must be produced, when it must be updated, or whether the record is sufficient for action.

Each artifact should perform a distinct function. The dependency disclosure package should identify the high-consequence dependency at issue, locate it within the relevant supplier pathway, and state why that relationship or node matters for continuity or integrity. The provenance record should capture the origin, custody, relevant processing steps, and any known substitutions or breaks in traceability. The change notification should report consequential rewiring events, specify what changed and when, and state the expected effect on continuity or integrity exposure. The resilience support annex should identify qualified alternatives, fallback paths, recovery constraints, and other assumptions that bear on whether the dependency can be absorbed, worked around, or restored in practice \parencite{ch5_dodi424515,ch5_dfars2522467007,ch5_nist800161r1upd1}. Regardless of final nomenclature, the purpose of these artifacts is to create a minimum evidentiary package that lets the Department determine what it knows, what changed, and what requires validation or action.

Post-award execution and sustainment require more than one-time disclosure because high-consequence dependencies do not remain fixed after contract award. Supplier substitutions, ownership changes, geographic relocation, obsolescence workarounds, qualification changes, and other forms of consequential rewiring can alter continuity or integrity exposure even when the end item being delivered appears unchanged \parencite{ch5_dodi424515,ch5_nist800161r1upd1}. A credible regime must therefore combine triggered reporting with periodic revalidation. Triggered reporting ensures that material changes are surfaced when they occur. Periodic revalidation ensures that high-consequence dependencies do not drift into stale assumptions simply because no single event was formally reported. The point is not continuous surveillance of every supplier transaction. It is to ensure that award-stage artifacts remain decision-useful over time and that the Department learns about consequential rewiring before it discovers the effects through schedule disruption, quality failure, or emergency substitution.

A visibility regime that stops at the prime contractor will reproduce the same ceiling it is meant to reduce. The Department usually reaches high-consequence sub-tier dependencies indirectly, which means visibility obligations must move through primes by way of flow-down requirements, recordkeeping duties, and prime-mediated reporting pathways rather than an unrealistic expectation of direct government access to every lower-tier firm \parencite{ch5_dfars2522467008,ch5_dfars2522397018,ch5_dodi424515}. The objective is not unrestricted disclosure of every supplier relationship. It is a structured way to ensure that when a designated dependency sits below Tier~1, the evidence needed for traceability, provenance, change reporting, and resilience assessment can still move upward to the point of oversight. That requires primes to do more than passively receive information. They must be responsible for flowing down designated visibility obligations, maintaining relevant records, and furnishing defined evidence to the government when high-consequence dependencies trigger review. Flow-down is not a secondary compliance detail. It is the institutional mechanism that gives the Department indirect reach into the parts of the supply chain where the visibility ceiling is often highest.

Requiring information is not the same as verifying it. A credible visibility regime must therefore treat verification as a risk-tiered function aligned to the Section~\ref{chap:network_analysis} workflow rather than as a uniform audit burden imposed across the entire supplier base. Dependencies that fall into the act now or validate then act categories should receive the deepest scrutiny, including documentary review, targeted supplier inquiry, and, where necessary, government access to supporting records or assessment mechanisms that test whether required controls are actually in place \parencite{ch5_dfars2522047020,ch5_dfars2522467007}. Dependencies in the monitor category may warrant lighter revalidation, exception-based review, or periodic refresh without the same investigative depth. The point is not to promise exhaustive verification. It is to concentrate scarce validation capacity where structural consequence and evidentiary uncertainty are highest, so the Department can distinguish between information that has merely been reported and information that is strong enough to support intervention.

The authorities needed to begin this shift are not wholly absent. DoD already has acquisition, program protection, sustainment, counterfeit avoidance, and pilot mapping channels through which a more disciplined visibility regime could be inserted, including DoDI 5000.85, the Technology and Program Protection Guidebook, DoDI 4245.15, DFARS 252.246-7007, DFARS 252.246-7008, and FY2024 NDAA Section 856 \parencite{ch5_dodi500085,ch5_dodi424515,ch5_dfars2522467007,ch5_ndaa2024_sec856}. The point of this section is therefore not to propose final clause text or statutory language. It is to identify the institutional insertion points through which actionable visibility can be made repeatable for high-consequence semiconductor dependencies. The logic is direct. The Department must decide where visibility enters the lifecycle, what evidence artifacts must exist, how obligations move beyond primes, and how verification is prioritized under constrained capacity. Once those levers are specified, the next question is how the resulting information should be governed across the enterprise, which is the subject of Section~\ref{sec:ch5_data_governance_models}.

%========================
% Section 5.5
%========================
\section{Data Governance Models}
\label{sec:ch5_data_governance_models}

Section~\ref{sec:ch5_acquisition_regulatory_levers} identified where actionable visibility should enter the acquisition lifecycle and what recurring artifacts should exist. This section takes up the next institutional question, which is how that evidence should be governed once it exists. This section compares three governance postures, centralized, federated, and hybrid, and argues that the Department should adopt a hybrid model that centralizes standards, oversight, and interoperability requirements without forcing all sensitive supplier data into a single repository. That question is not administrative detail. Once dependency disclosure packages, provenance evidence records, change notifications, and resilience support annexes begin to exist, the Department must decide how those records will be stored, exchanged, queried, protected, and used across programs and primes. A visibility regime that lacks a workable governance model will either fragment into contractor-specific silos or harden into a brittle system that firms resist, and the Department cannot safely manage \parencite{ch5_nist800161r1upd1,ch5_dbb2025_illumination,U.S.GovernmentAccountabilityOffice2025}.

Data governance is not a back-end technical problem. It determines whether the visibility regime can satisfy the principles set out in Section~\ref{sec:ch5_policy_objectives_design_principles} once sensitive evidence begins to circulate. If dependency records remain scattered across contractor portals, program offices, and prime-specific systems, the Department cannot compare artifacts across firms, detect shared sub-tier vulnerabilities, or see how risk accumulates across programs. If the same evidence is centralized too aggressively, the government creates a high-value cyber target and a disclosure regime that firms may treat as incompatible with minimum necessary disclosure and legitimate proprietary protection. Governance is therefore where interoperability, verifiability, selective disclosure, and burden-aware implementation become operational rather than rhetorical. The question is not simply where records live. It is how the Department can obtain usable cross-program insight without recreating either the siloed visibility ceiling or a centralized data honeypot of its own making \parencite{ch5_nist800161r1upd1,ch5_dbb2025_illumination,U.S.GovernmentAccountabilityOffice2025,ch5_farsubpart423,ch5_dscsa_interop_2023}.

A centralized governance model has real advantages and should be taken seriously on its own terms. If the Department were to require that visibility artifacts be submitted into a common repository governed by shared standards and oversight rules, it would gain a clearer enterprise-wide picture of supplier dependencies, simpler cross-program querying, more consistent auditability, and a stronger basis for identifying shared vulnerabilities that no single prime or program office can see alone \parencite{ch5_nist800161r1upd1,ch5_dbb2025_illumination,U.S.GovernmentAccountabilityOffice2025,ch5_farsubpart423}. Those are not trivial benefits. They speak directly to the current fragmentation problem. But centralization alone is still not a satisfactory answer. A single repository that aggregates detailed dependency, provenance, and resilience evidence across programs would create an attractive cyber and insider target, heighten contractor concern about proprietary exposure, and pressure the Department toward broader data collection than decision use actually requires. A fully centralized model therefore advances interoperability and oversight, but it does so by creating new risks that sit uneasily with minimum necessary disclosure, selective access, and burden-aware implementation.

A federated or permissioned model also has clear strengths and, in some respects, fits the section's principles more naturally. If sensitive supplier evidence remains under contractor or distributed control, firms can disclose only what decision use requires, protect proprietary relationships more effectively, and limit the cyber and competitive risks associated with broad centralized exposure \parencite{ch5_nist800161r1upd1,ch5_dscsa_interop_2023,ch5_ntia_sbom_2021,ch5_nist_ir_8419}. That posture is more consistent with minimum necessary disclosure and may lower resistance to participation. But a federated model is not sufficient simply because it is less intrusive. If each prime, program, or vendor maintains its own closed visibility environment, the Department still cannot compare artifacts across firms, detect shared sub-tier chokepoints, or see how risk accumulates across the enterprise. A Raytheon solution, a Boeing solution, and an Anduril solution do not add up to usable governance visibility for the Department. Without governed interoperability, a federated model becomes another name for fragmentation.

The Department should therefore adopt a hybrid governance model. Standards, minimum data fields, identifiers, exchange protocols, audit logic, and access rules should be governed centrally so that visibility artifacts are comparable across programs and primes and usable for enterprise-level oversight. But underlying records, sensitive supplier details, and some higher-resolution provenance evidence should remain distributed or permissioned where proprietary, export-control, classification, or cyber-risk concerns justify tighter control \parencite{ch5_nist800161r1upd1,ch5_dbb2025_illumination,ch5_farsubpart423,ch5_dscsa_interop_2023,ch5_nist_ir_8419}. The governing principle is not to centralize every byte. It is to centralize the rules that make evidence interoperable and reviewable while decentralizing storage and disclosure to the degree consistent with decision use. That approach best fits the section's design principles because it preserves minimum necessary disclosure, supports cross-program visibility, and reduces the risk that the Department tries to solve fragmentation by building a single repository that industry will resist or adversaries will target.

A hybrid governance posture only works if the resulting regime is genuinely interoperable. The Department does not need one monolithic platform, but it does need one governed ecosystem in which visibility artifacts can be exchanged, interpreted, and validated across primes, programs, and tiers. That requires more than broad compatibility claims. It requires common identifiers, machine-readable minimum fields, shared event logic for changes such as substitution or geography shifts, standardized exchange protocols, and rules for how records are authenticated, updated, and queried across institutional boundaries \parencite{ch5_nist800161r1upd1,ch5_dscsa_interop_2023,ch5_ntia_sbom_2021,ch5_nist_ir_8536_2,ch5_nist_ir_8419}. Interoperability is therefore not a technical convenience layered on top of governance. It is the condition that prevents a hybrid model from collapsing back into disconnected contractor portals and program-specific silos. Without it, the Department may receive more records, but it will still struggle to see the cross-prime and sub-tier vulnerabilities that matter most.

Good governance is not only a question of where data resides. It is also a question of who can see what, under what conditions, and with what accountability. The Department should therefore govern these visibility artifacts through least-privilege and role-based access rules that distinguish among program execution, enterprise oversight, supplier validation, and broader analytic use. Sensitive supplier evidence should be disclosed selectively rather than universally, with access segmented by program relevance, risk designation, and information sensitivity, and with auditable logs, revocation procedures, and retention rules that make record use traceable over time \parencite{ch5_nist800161r1upd1,ch5_dscsa_interop_2023,ch5_nist_ir_8536_2,ch5_nist_ir_8419,ch5_farsubpart423}. That same governance logic should extend to classification and handling boundaries. Some records may remain unclassified but commercially sensitive, others may require controlled access because they expose sourcing strategies, process details, or vulnerable chokepoints, and still others may need stronger protection when they intersect with classified programs or mission-critical functions. A workable regime therefore does not promise universal transparency. It creates a controlled disclosure architecture in which the Department can obtain the evidence necessary for assurance without demanding indiscriminate access to everything firms know.

One practical model would be a governed interoperability ecosystem operated by approved infrastructure providers under Department-defined standards, access rules, and audit requirements. In such a model, firms would not be required to surrender all supplier data into a single government-owned repository. Instead, each participating firm would retain control over its sensitive underlying records while issuing standardized, machine-readable visibility artifacts that document designated dependency, provenance, and change events in a form that can be authenticated and linked across tiers. Higher-resolution supplier data would remain permissioned. Prime contractors would furnish the minimum evidence required for decision use and, when validation is triggered, coordinate selective disclosure from relevant lower-tier firms through the same governed framework. The result would not be universal visibility into every commercial relationship. It would be a structured chain of evidence that allows authorized users to trace dependencies, assess provenance claims, detect consequential rewiring, and compare signals across primes and programs without forcing every proprietary detail into one place \parencite{ch5_dscsa_interop_2023,ch5_ntia_sbom_2021,ch5_nist_ir_8536_2,ch5_space_trustworthy_data_2025}.\footnote{McGillis and Reed (2026) is cited as a conference paper presented at the 12th IAA Space Traffic Management Conference (February 2026). The substantive claims are independently supported by the published works of Reed and colleagues, including NIST IR 8536 (Pease et al., 2025) and NIST IR 8419 (Stouffer et al., 2022).}

Such an ecosystem could support more than passive record exchange. Specialized third-party providers could furnish secure identity and credentialing services, exchange infrastructure, verification tools, audit logging, and analytic services that help the Department identify shared chokepoints, anomalous substitutions, or emerging cross-prime vulnerabilities. AI-enabled analysis could then operate on a stronger evidentiary foundation because the underlying records would be structured, auditable, and governed rather than merely self-reported or manually assembled. The preceding illustration is not a recommendation that the Department adopt any one protocol, ledger, vendor, or platform. It shows what the section's policy principles could make possible. The recommendation is functional rather than technological. Any acceptable implementation must support interoperable exchange, selective disclosure, auditable provenance, governed third-party participation, and enterprise-level analysis without creating a single centralized data honeypot \parencite{ch5_nist800161r1upd1,ch5_space_trustworthy_data_2025}.

A visibility regime without a sound governance model will fail in one of two ways. It will either fragment into prime-specific systems that never produce Department-level insight, or it will over-centralize into a brittle repository that firms resist and adversaries would prize. The appropriate posture is therefore a hybrid governance ecosystem that centralizes standards, interoperability rules, oversight, and audit logic while keeping underlying sensitive records distributed or permissioned to the degree consistent with decision use. That posture is the best fit with the section's principles because it preserves minimum necessary disclosure, supports verifiability, makes interoperability real, and gives the Department a way to see cross-prime and sub-tier vulnerabilities without demanding indiscriminate access to everything firms know. Once that governance posture is established, however, another problem becomes unavoidable. A workable regime still depends on whether firms have incentives to participate, whether obligations are enforceable, and whether the burden of compliance is allocated in a way that smaller suppliers can sustain. The next section therefore turns to incentives, enforcement, and cost-sharing.

%========================
% Section 5.6
%========================
\section{Incentives, Enforcement, and Cost-Sharing}
\label{sec:ch5_incentives_enforcement_cost_sharing}

Section~\ref{sec:ch5_data_governance_models} explained how visibility artifacts should be governed once they exist. But acquisition levers and data governance are still not enough to make the regime durable in practice. A visibility system only works if firms have reasons to participate, if obligations are enforceable when they do not, and if the burden of compliance is distributed in a way that lower-tier suppliers can actually sustain. That is especially important in the semiconductor supply chain, where proprietary relationships and process knowledge create strong incentives to withhold information even as hidden common vulnerabilities can expose multiple firms at once. The policy problem in this section is therefore not only how to require evidence, but how to make participation rational, noncompliance consequential, and mitigation feasible when the weakest firms in the chain may also be among the most consequential \parencite{ch5_nist800161r1upd1,ch5_dod2022_supplychains,U.S.GovernmentAccountabilityOffice2025}.

The Department should therefore adopt a mixed compliance model. A regime built on voluntary participation alone will not overcome the proprietary incentives and information asymmetries that sustain the visibility ceiling. A regime built only on punitive enforcement will encourage defensive reporting, concealment, and formal compliance without useful disclosure. The workable alternative is to combine limited but meaningful incentives, real consequences for noncompliance, safe harbor for timely disclosure and corrective action, and scale-appropriate obligations that reflect both the consequence of the dependency and the capacity of the firm \parencite{ch5_nist800161r1upd1,ch5_dod2022_supplychains,U.S.GovernmentAccountabilityOffice2025}. That approach treats the visibility requirement not as a moral appeal, but as a governed operating requirement that firms can rationally participate in and the Department can credibly enforce.

Incentives matter because the visibility problem is partly a collective-action problem. Firms have strong reasons to protect supplier relationships, process knowledge, and sourcing strategies, but that same opacity can conceal lower-tier chokepoints and common dependencies that expose multiple firms simultaneously to disruption, delay, or substitution risk. In that sense, better participation in this visibility regime is not only a compliance burden. It can also reduce shared exposure to vulnerabilities that no one firm can fully see or mitigate on its own \parencite{ch5_dod2022_supplychains,U.S.GovernmentAccountabilityOffice2025,ch5_dbb2025_illumination}. The Department therefore has reason to attach concrete advantages to timely and reliable participation, especially for firms that can furnish auditable visibility artifacts, disclose consequential changes early, and support validation when asked. Those advantages could include reduced review friction when reporting is strong and timely, access to pilot efforts or technical support, and clearer treatment of certain compliance costs where the Department is imposing sustained assurance requirements. Incentives are not a substitute for enforcement. They are a way to make truthful participation, early disclosure, and sustained recordkeeping more attractive than opacity, delay, or minimal compliance.

Enforcement must be real because a visibility regime with no meaningful consequence for noncompliance will be treated as advisory. The Department does not need a maximalist penalty structure. It needs to make clear that supplier visibility is part of program accountability rather than an optional reporting preference. Where firms fail to furnish required artifacts, withhold material changes, or do not support validation when triggered, the consequences should affect contract performance management, review intensity, corrective-action obligations, and, for designated high-consequence work, continued eligibility to perform without heightened oversight \parencite{ch5_dfars2522467007,ch5_dfars2522047020}. That posture is already consistent with how the Department treats other assurance-related requirements, including counterfeit avoidance systems and assessment-based verification of cybersecurity controls \parencite{ch5_dfars2522467007,ch5_dfars2522047020}. The point is not to punish every reporting defect as if it were fraud. It is to ensure that firms understand noncompliance as a real operational and contractual risk, not a paperwork lapse that can be ignored until a disruption exposes it \parencite{ch5_nist800161r1upd1}.

Safe-harbor logic should distinguish between firms that surface problems early and firms that conceal them until performance fails. The Department should give more favorable treatment to contractors that disclose traceability breaks, provenance uncertainty, supplier substitutions, or other consequential changes promptly, preserve the relevant records, and take corrective action in good faith. That approach already has precedent. DoD cost rules for counterfeit and suspect counterfeit electronic parts tie the allowability of certain costs to the existence of an approved detection and avoidance system, use of appropriate sourcing channels, and timely written notice after the contractor becomes aware of the problem \parencite{ch5_dfars23120571}. Counterfeit avoidance policy likewise assumes that detection, reporting, quarantine, and corrective response matter more than silent noncompliance \parencite{ch5_dfars2522467007}. Safe harbor should therefore reduce the penalty for timely disclosure and responsible correction, not eliminate accountability altogether. Firms that conceal material changes, delay notice, or repeatedly fail to maintain required evidence should not receive the same treatment as firms that surface problems while mitigation is still possible.

Compliance is not costless. Producing visibility artifacts, maintaining traceability and provenance records, supporting selective disclosure, and responding to validation requests all require labor, systems, and management attention over time. The Department therefore cannot treat cost as an afterthought once it decides to impose sustained assurance requirements. It must distinguish between obligations firms should absorb as part of ordinary performance and obligations whose cost should be treated as allowable or suitable for targeted cost-sharing when the Department is asking industry to maintain a higher level of evidence than current practice normally requires. DoD already has precedent for linking cost treatment to compliance behavior in the counterfeit-parts context, where allowability turns in part on the existence of an approved system and timely notice after a problem is discovered \parencite{ch5_dfars23120571}. The same logic can be extended here. Where visibility and mitigation requirements address systemic lower-tier vulnerabilities that no one firm can fully solve on its own, some compliance costs should be understood not only as firm-specific burdens but as resilience investments with broader industrial-base benefit \parencite{ch5_dod2022_supplychains,ch5_nist800161r1upd1}. The Department should not assume that the weakest firms in the chain can absorb those costs alone simply because they sit furthest from the prime.

Scale-appropriate obligations are essential because the firms least able to absorb new compliance burdens may also sit at some of the most consequential points in the supply chain. Section~\ref{sec:ch5_policy_objectives_design_principles} already argued that burden-aware implementation is not a concession to opacity. It is a condition of durable participation below Tier~1. The implication is not that small or specialized firms should be exempt from visibility requirements simply because they are small. It is that the depth, frequency, and evidentiary burden of those requirements should be calibrated to consequence, role in the chain, and practical capacity rather than imposed uniformly across all actors. Federal traceability regimes in other sectors already show that meaningful record-based compliance can be structured around defined events and defined data elements without requiring identical burdens for every participant \parencite{ch5_usc2223,ch5_fda_small_entity_traceability_guide_2025}. For the Department, that means preserving the same assurance objective for high-consequence dependencies while allowing the compliance pathway to vary by scale, provided the resulting records remain usable for governance. A small firm should not be asked to build the same infrastructure as a prime contractor. It should be asked to furnish the evidence necessary for the role it plays in a way it can actually sustain.

These elements are most effective when they operate as a single compliance model rather than as separate policy tools. Incentives make early and truthful participation more attractive. Enforcement makes clear that visibility obligations are part of real program accountability. Safe harbor distinguishes good-faith disclosure from concealment. Cost treatment recognizes that some assurance burdens create benefits that extend beyond the firm that bears them first. Scale-appropriate obligations make it possible to carry the regime below Tier~1 without hollowing it out through uniform demands. This is what a durable visibility regime requires. It must make participation rational, noncompliance consequential, corrective disclosure worthwhile, and lower-tier adoption feasible enough to sustain over time \parencite{ch5_nist800161r1upd1,ch5_dod2022_supplychains,U.S.GovernmentAccountabilityOffice2025}.

The next question is how to phase this regime into practice. The Department should not attempt to impose it everywhere at once. It should phase requirements in through designated semiconductor-relevant categories, beginning where Section~\ref{chap:network_analysis} identifies the strongest act now and validate then act priorities and then widening as reporting, governance, and validation capacity mature. A finite rollout matters for two reasons: it allows the Department to learn where artifacts, access rules, and verification pathways work as intended, and it prevents the regime from collapsing under its own ambition before lower-tier participation is sustainable. Section~\ref{sec:ch5_implementation_roadmap} therefore sets out an implementation roadmap that sequences pilots, scaling, and institutionalization against the act now, validate then act, and monitor workflow rather than assuming enterprise-wide adoption from the start \parencite{ch5_ndaa2024_sec856,ch5_nist800161r1upd1,ch5_dod2022_supplychains}.

%========================
% Section 5.7
%========================
\section{Implementation Roadmap}
\label{sec:ch5_implementation_roadmap}

Section~\ref{sec:ch5_acquisition_regulatory_levers} defined what evidence the Department should require, Section~\ref{sec:ch5_data_governance_models} explained how that evidence should be governed, and Section~\ref{sec:ch5_incentives_enforcement_cost_sharing} addressed how participation should be sustained and enforced. What remains is not whether the regime is justified, but how to phase it into practice without trying to impose it everywhere at once. This section therefore sets out a bounded implementation roadmap that begins with pilots, uses consequence rather than convenience to determine where the Department starts, and expands only as the underlying artifacts, governance rules, and validation pathways are shown to work in practice \parencite{ch5_nist800161r1upd1,ch5_dbb2025_illumination}.

The Department does not need to wait for new legislation to begin. Section 856 of the FY2024 National Defense Authorization Act already directs the Under Secretary of Defense for Acquisition and Sustainment to establish a pilot program to analyze, map, and monitor supply chains for up to five covered weapons platforms, employing a combination of government and commercial tools for that purpose \parencite{ch5_ndaa2024_sec856}. The accompanying House Armed Services Committee Report (H.Rept.\ 118-125) further directs the Department to assess the feasibility of a continuously monitored system leveraging artificial intelligence to identify and mitigate supply chain risks. GAO reports that DoD is already moving in this direction through the Supply Chain Risk Evaluation Environment (SCREEn) and identifies the five planned systems as the F-35, MV-22C Osprey, Patriot, C-5C/M Super Galaxy, and Standard Missile 6 \parencite{U.S.GovernmentAccountabilityOffice2025}. That existing authority matters because it makes the roadmap in this section immediate rather than hypothetical. The question is not whether the Department can begin piloting now, but whether those pilots will be designed narrowly as platform exercises or strategically as entry points into the broader semiconductor dependency network.

Covered-platform pilots are a sensible starting point because they provide bounded scope, clear program ownership, and an authority structure the Department can use immediately. A platform-centered pilot makes it easier to assign responsibility, define the initial reporting boundary, and test whether artifacts, access rules, and validation procedures work under real acquisition conditions. But platform pilots are not sufficient if they illuminate only one weapon system at a time. As Sections~\ref{chap:dod-supply-chain} and~\ref{chap:network_analysis} show, the highest-consequence vulnerabilities in the DoD-linked semiconductor network are often concentrated in shared upstream pathways, chokepoints, bridges, and corridor intermediaries that are not visible from any single prime-facing view alone. The policy value of a pilot therefore lies not only in what it reveals about one covered platform, but in whether it exposes the broader semiconductor dependencies and recurring validation problems that cut across programs. DoD's own supply-chain strategy points in the same direction by emphasizing cross-program analysis, common standards, and industrial-base learning rather than isolated program-level visibility alone \parencite{ch5_dod2022_supplychains,U.S.GovernmentAccountabilityOffice2025}.

The governing rule for this roadmap is simple: platforms are the entry point, but the object of learning is the broader semiconductor dependency network that sits behind them. Each pilot should therefore be designed to answer two questions at once. It should identify which dependencies are consequential for the covered platform itself, and it should reveal which upstream suppliers, chokepoints, corridor intermediaries, provenance gaps, or substitution problems recur across multiple programs once the evidence is structured and compared. That is the point at which a bounded platform exercise begins to generate Department-level value. A pilot that treats the platform as the final analytic boundary will therefore miss part of what the Department most needs to learn. For that reason, the roadmap proceeds in three phases: an initial bounded pilot to prove the core regime in practice, a second phase that expands cross-program comparison and governed interoperability, and a final phase that institutionalizes the regime into routine acquisition and oversight for designated semiconductor-relevant categories.

The first phase should be a bounded pilot focused on designated semiconductor-relevant dependencies within or adjacent to the covered platforms authorized under Section 856 \parencite{ch5_ndaa2024_sec856}. Its purpose is not to illuminate the entire supply network at once. It is to prove that the core regime can actually function in practice for a limited set of high-consequence dependencies. In this phase, the Department should test whether dependency disclosure packages, provenance evidence records, change notifications, and resilience support annexes can be produced in usable form, governed under the hybrid model proposed in Section~\ref{sec:ch5_data_governance_models}, and selectively validated when the Section~\ref{chap:network_analysis} workflow indicates that uncertainty is still decision-relevant. A successful Phase I pilot would therefore do more than show that records can be collected. It would show that the Department can obtain auditable evidence, protect sensitive data, trigger additional disclosure when required, and use the resulting information to distinguish between stable priorities and cases that still require targeted validation, which is exactly the kind of pilot-to-practice transition the Defense Business Board argues is necessary \parencite{ch5_dbb2025_illumination}.

The second phase should begin only after the initial pilot demonstrates that the core artifacts are usable, and the governance controls hold under real conditions. Its purpose is to move beyond one platform at a time and test whether the regime can generate cross-program visibility into shared semiconductor dependencies. In this phase, the Department should compare artifacts across covered systems or dependency clusters, identify recurring lower-tier suppliers and chokepoints, and test whether governed interoperability allows evidence to be queried and validated across primes without collapsing back into isolated contractor portals. That shift is consistent with DoD's own supply-chain strategy, which calls for common standards and cross-program analysis rather than isolated program-level visibility \parencite{ch5_dod2022_supplychains}. A successful second phase would therefore show not only that one platform can see its dependencies, but that the Department can begin to detect shared lower-tier vulnerabilities that no single program can identify alone.

The third phase should begin only after the Department has shown that the artifacts are usable, the governance controls hold, and cross-program comparison produces insight that matters for action. At that point, the regime should move from bounded pilot activity into routine acquisition and oversight practice for designated semiconductor-relevant categories. That means the core artifacts should become standing program requirements where consequence justifies them, the hybrid governance rules should become part of normal evidence handling rather than a pilot exception, and validation should be integrated into ordinary oversight and revalidation cycles rather than treated as a temporary experiment. This is the point at which phased implementation becomes durable institutional practice, with enduring roles, reporting, reassessment, and enterprise use rather than one-time deployment or disconnected pilots \parencite{ch5_nist800161r1upd1,ch5_dbb2025_illumination}. A successful third phase therefore marks the point at which semiconductor visibility stops being a special project and becomes a normal part of how the Department manages high-consequence dependencies.

Expansion should be conditional rather than automatic. The Department should move from one phase to the next only if the pilot demonstrates that the core artifacts are usable for decision-making, that consequential changes are reported with enough timeliness to matter, that selective disclosure and access controls work without breaking the flow of evidence, and that validation can distinguish between stable findings and cases that still require targeted inquiry. Expansion should also depend on whether the hybrid governance model is maturing into a functioning ecosystem rather than remaining trapped inside one program office or one contractor stack. In practical terms, that means the Department should test whether approved third-party infrastructure or equivalent exchange mechanisms can support interoperable record sharing, audit logging, governed queries, and selective disclosure across firms and programs without forcing all sensitive data into one place \parencite{ch5_nist_ir_8536_2,ch5_nist_ir_8419}. NIST's implementation guidance is useful here because it treats rollout as a phased roadmap tied to roles, reporting, reassessment, and adjustment rather than one-time deployment \parencite{ch5_nist800161r1upd1}. The Department should also be willing to pause or redesign a phase if participation falls off sharply below Tier~1, if artifact quality is too weak for action, if the ecosystem remains siloed by prime or program, or if the pilot generates records without yielding clearer decisions about where to act, validate, or monitor. A pilot is not successful merely because it produces dashboards or more data. It is successful if it reduces uncertainty where consequence is highest and creates a basis for broader institutionalization rather than another isolated analytic effort \parencite{U.S.GovernmentAccountabilityOffice2025,ch5_dbb2025_illumination}.

The roadmap proposed here is therefore not a call for enterprise-wide rollout by fiat. It is a bounded sequence that begins under the Section 856 pilot authority, uses covered platforms as practical entry points, and expands only if the Department can produce usable artifacts, govern them through shared standards and access-controlled data exchange, and validate high-consequence findings at a sustainable pace \parencite{ch5_ndaa2024_sec856,ch5_nist800161r1upd1}. If those conditions are met, the pilots become more than isolated demonstrations. They become the basis for a broader semiconductor visibility regime that can move from pilot mapping to an enduring assurance regime. GAO and the Defense Business Board both point to the same underlying challenge, which is not whether DoD can run pilots, but whether it can turn them into durable practice rather than another fragmented analytic effort \parencite{U.S.GovernmentAccountabilityOffice2025,ch5_dbb2025_illumination}. The final task of this section is therefore to state what the Department and Congress should do now in concrete terms. Section~\ref{sec:ch5_recommendations_dod_congress} turns to those recommendations directly.

%========================
% Section 5.8
%========================
\section{Policy Recommendations}
\label{sec:ch5_recommendations_dod_congress}

Sections~\ref{sec:ch5_acquisition_regulatory_levers} through~\ref{sec:ch5_implementation_roadmap} developed the operating logic of this section. They specified what evidence the Department should require, how that evidence should be governed, how participation and enforcement should be structured, and how implementation should be phased. This section now states the section's recommendations in direct terms. The aim is not to restate the full argument, but to identify the highest-value actions that the Department of Defense and Congress can take now to build a defensible semiconductor visibility and validation regime. Because this section is intended to be usable by decision-makers, the recommendations are stated as concrete actions rather than as general policy aspirations.

The Department of Defense can begin under existing authority. The Office of the Under Secretary of Defense for Acquisition and Sustainment (OUSD(A\&S)) should use the Section 856 pilots not only to produce platform-bounded dependency maps, but also to generate a bounded cross-platform microelectronics learning product that identifies shared lower-tier suppliers, recurring chokepoints, and repeated validation problems across the covered systems \parencite{ch5_ndaa2024_sec856,U.S.GovernmentAccountabilityOffice2025}. At the same time, the Department should define and require a small standardized set of audit-ready visibility artifacts for high-consequence dependencies and stand up a risk-tiering process that determines which dependencies warrant the deepest visibility and validation burden. Those steps should not wait for a wholly new statutory framework because they are the minimum conditions for turning existing pilot activity into a disciplined operating model rather than another round of ad hoc supply-chain illumination \parencite{ch5_dod2022_supplychains,ch5_nist800161r1upd1}.

To make the regime durable, the Department should institutionalize the operating conditions that make visibility usable rather than episodic. OUSD(A\&S), working with the Chief Digital and Artificial Intelligence Office (CDAO), should establish one common standards framework for identifiers, minimum fields, exchange rules, audit logs, and access controls so that visibility artifacts can be compared and queried across programs without manual reconciliation \parencite{ch5_nist800161r1upd1,ch5_dbb2025_illumination}. The Department should govern sensitive supplier evidence through a permissioned model rather than a single centralized repository, require event-triggered updates and periodic revalidation for designated high-consequence dependencies, and adopt the mixed compliance model developed in Section~\ref{sec:ch5_incentives_enforcement_cost_sharing} so that disclosure, validation, and lower-tier participation remain sustainable over time \parencite{ch5_dod2022_supplychains,ch5_nist800161r1upd1}. It should also build an approved third-party ecosystem for exchange, validation, audit, and analytics under Department-defined standards, because a visibility regime that depends only on isolated program offices or prime-specific systems will not scale into enterprise practice \parencite{ch5_nist_ir_8419,ch5_dbb2025_illumination}.

Congress should now reinforce what the Department has begun. Section 856 is a useful starting point, but it remains a bounded pilot authority tied to up to five covered weapons platforms rather than a durable framework for semiconductor visibility and validation across the defense enterprise \parencite{ch5_ndaa2024_sec856}. Congress should therefore extend and broaden that authority into a standing semiconductor visibility and validation program that preserves the value of bounded pilots while making clear that the long-term objective is institutionalized cross-platform learning, not a temporary demonstration. That change matters because the Department's challenge is no longer whether it can generate interesting pilot outputs. It is whether it can convert those pilots into a durable operating model for identifying shared lower-tier dependencies, common chokepoints, and recurring validation problems before they manifest as readiness or production failures \parencite{U.S.GovernmentAccountabilityOffice2025,ch5_dbb2025_illumination}.

Congress should also resource the regime in a way that matches the scale of the problem. A standing semiconductor visibility and validation program will not become durable if DoD is expected to sustain it through episodic studies, borrowed personnel, or ad hoc analytic contracts. Congress should appropriate resources for validation capacity, governance infrastructure, standards development, approved third-party ecosystem support, and targeted mitigation where systemic lower-tier vulnerabilities are identified \parencite{ch5_dod2022_supplychains,ch5_nist800161r1upd1}. That includes the capacity to review artifacts, govern access and exchange rules, support interoperable implementation, and help address vulnerabilities in smaller but structurally important firms that may be the least able to absorb the cost of mitigation on their own.

Table~\ref{tab:ch5_action_matrix} translates these recommendations into an actor-specific action matrix for semiconductor visibility and validation. Its purpose is not to restate the section in compressed form. It is to identify, in one place, what should change, who should own the change, what implementation hook can carry it, and what success should look like if the recommendation is taken seriously. Read this table as an action tool rather than a summary table.  The recommendations show how the Department and Congress can move from fragmented visibility efforts toward a defensible semiconductor assurance regime grounded in clear ownership, concrete implementation hooks, and observable signals of progress.

These recommendations form a near-term agenda for moving from fragmented illumination efforts to a decision-grade semiconductor visibility and validation regime. They do not ask the Department to begin from scratch, and they do not depend on a single platform, vendor, or statutory overhaul before progress can begin. Instead, they identify the specific actions through which existing authority can be used more coherently, the points at which congressional reinforcement becomes necessary, and the conditions under which pilots can mature into durable practice. The final section of this section returns to the broader argument and clarifies what it means, in policy terms, to move from illumination to assurance.

% --------------- Chapter 5 Action Matrix --------------
\begin{landscape}
\pagestyle{landscape}
\begin{table}[!p]
  \centering
  \footnotesize
  \renewcommand{\arraystretch}{1.1}
  \setlength{\tabcolsep}{4pt}
  \captionsetup{skip=4pt}

  \caption[\textit{Action Matrix for DoD and Congressional Recommendations}]
{\textit{Action Matrix for DoD and Congressional Recommendations}}

  \label{tab:ch5_action_matrix}

  % chktex-file 8   en-dash compounds are correct academic usage
% chktex-file 36  acronym parentheses (OUSD(A&S)) need no leading space
\begin{tabular}{@{}>{\raggedright\arraybackslash}p{8.6cm}
                >{\raggedright\arraybackslash}p{1.9cm}
                >{\raggedright\arraybackslash}p{4.5cm}
                >{\raggedright\arraybackslash}p{7.0cm}@{}}

\toprule
\textbf{Recommendation} & \textbf{Owner} & \textbf{Implementation hook} & \textbf{Success signal} \\
\midrule

\textbf{Use the Section 856 pilots to produce platform-bounded dependency maps and a bounded cross-platform semiconductor learning product.}
& OUSD(A\&S)
& Section 856 pilot guidance; A\&S memorandum; annual pilot reporting
& Each pilot produces a program-level dependency map and a cross-platform semiconductor view identifying lower-tier suppliers and chokepoints \\
\specialrule{0.3pt}{0pt}{0pt}

\textbf{Define and require a standardized set of decision-grade visibility artifacts for high-consequence dependencies.}
& OUSD(A\&S); DPC
& A\&S pilot guidance; acquisition policy; later DoDI/DFARS updates
& Programs generate machine-readable disclosure packages, provenance records, change notices, and resilience annexes \\
\specialrule{0.3pt}{0pt}{0pt}

\textbf{Establish a risk-tiering process to determine which dependencies receive deeper visibility and validation requirements.}
& OUSD(A\&S); CAEs/PEOs
& A\&S designation criteria; pilot guidance; program protection reviews
& DoD identifies high-consequence dependencies and assigns appropriate oversight or monitoring \\
\specialrule{0.3pt}{0pt}{0pt}

\textbf{Create a common standards framework for identifiers, minimum fields, exchange rules, audit logs, and access controls.}
& OUSD(A\&S); DoD CDAO
& A\&S--CDAO standards memorandum; pilot data schema; later DoD policy
& Visibility artifacts can be exchanged and queried across programs without manual reconciliation \\
\specialrule{0.3pt}{0pt}{0pt}

\textbf{Govern sensitive supplier evidence through a permissioned model rather than a single centralized repository.}
& OUSD(A\&S); DoD CDAO
& Pilot governance rules; access-control policy; departmental data guidance
& DoD validates supplier evidence while sensitive data remains selectively disclosed and auditable \\
\specialrule{0.3pt}{0pt}{0pt}

\textbf{Require event-triggered updates and periodic revalidation for designated high-consequence dependencies.}
& OUSD(A\&S); CAEs/PEOs
& Pilot reporting rules; acquisition guidance; DoDI 4245.15 / DFARS updates
& Substitution, ownership, geography, or certification changes are surfaced in time to inform decisions \\
\specialrule{0.3pt}{0pt}{0pt}

\textbf{Adopt a mixed compliance model combining incentives, enforcement, safe harbor for timely disclosure, and scale-appropriate obligations.}
& OUSD(A\&S); DPC
& Pilot compliance guidance; contract policy memoranda; later DFARS update
& Firms disclose problems earlier and noncompliance produces meaningful oversight consequences \\
\specialrule{0.3pt}{0pt}{0pt}

\textbf{Build an approved third-party ecosystem for exchange, validation, audit, and analytics under Department-defined standards.}
& OUSD(A\&S); DoD CDAO
& Pilot ecosystem criteria; approved provider framework; governance rules
& Third parties support exchange, audit logging, governed queries, and analytics without vendor lock-in or centralized data risks \\
\specialrule{0.3pt}{0pt}{0pt}

\textbf{Extend and broaden pilot authority into a standing semiconductor visibility and validation program.}
& Congress
& NDAA authorization or report language
& DoD moves from temporary pilots to a durable statutory framework for semiconductor visibility \\
\specialrule{0.3pt}{0pt}{0pt}

\textbf{Appropriate resources for validation capacity, governance infrastructure, standards development, and lower-tier mitigation.}
& Congress
& Appropriations language; NDAA direction; program elements
& DoD sustains validation capacity and targeted support for structurally important lower-tier suppliers \\

\bottomrule
\end{tabular}

\end{table}
\end{landscape}

%========================
% Section 5.9
%========================
\section{Synthesis: From Illumination to Assurance}
\label{sec:ch5_synthesis}

Section~\ref{chap:from_illumination_to_assurance} translated the visibility ceiling identified in Sections~\ref{chap:intro} through~\ref{chap:network_analysis} into an institutional response. It defined decision-grade visibility as the relevant policy objective, established the principles that should govern any serious attempt to reduce opacity, distinguished among traceability, provenance, resilience, and decision integration as separate assurance functions, and then translated that framework into acquisition levers, governance choices, compliance logic, a phased implementation roadmap, and concrete recommendations for DoD and Congress. The section's output is therefore not a generic call for more supply-chain transparency. It specifies what it would mean, in policy terms, to move from episodic illumination toward a selective, auditable, and action-oriented assurance regime for high-consequence semiconductor dependencies.

Section~\ref{chap:from_illumination_to_assurance}'s policy design did not stand apart from this report's earlier findings. It followed directly from the empirical and methodological argument developed in Sections~\ref{chap:dod-supply-chain} and~\ref{chap:network_analysis}, which showed that high-consequence semiconductor dependencies were structurally concentrated, often partially hidden, and not reliably manageable through prime-facing visibility alone. That finding shifted the policy problem from generic supply-chain awareness to targeted assurance under partial observability. The recommendations therefore did not rest on an abstract preference for more disclosure. They rested on this report's demonstration that dependencies most likely to matter for readiness and continuity could sit beyond the view of any single program or contractor, and that decision-making consequently required a more disciplined way to identify, verify, and monitor them.

The policy standard advanced in this section was not omniscience or universal transparency. It was decision-grade visibility: a selective, risk-tiered, and auditable regime that could produce enough evidence to identify, verify, monitor, and act on the dependencies that mattered most. That standard was deliberately bounded because the analysis did not eliminate the underlying limits of the problem. Some dependencies remained partially hidden, some relationships remained difficult to verify with complete confidence, and some integrity questions still exceeded what a provenance-aware but incomplete evidence base could certify. Those residual blind spots do not negate the value of the section's policy design. They define its interpretive boundaries and frame the Section~\ref{chap:conclusion} discussion of remaining limits, unresolved questions, and next steps.

Section~\ref{chap:conclusion} now steps back from the institutional design developed here and returns to the report as a whole. It restates the strategic claim that semiconductors have become a center-of-gravity dependency for the Department of Defense, summarizes the report's empirical, methodological, and policy findings, and clarifies what the network-based approach developed in this report changes for decision-making under partial observability. It also addresses the remaining limits of the evidence, the boundaries of the inferences drawn here, and the most important next steps for research and institutionalization. If Section~\ref{chap:from_illumination_to_assurance} showed how the Department could begin to build assurance in practice, Section~\ref{chap:conclusion} explains what that policy agenda means for this report's broader argument.